\section{Evolutionary Code and Algorithm Discovery}
\label{sec:evolutionary-code}

The self-improvement methods discussed in Section~\ref{sec:self-improvement} primarily operate within the space of model outputs---improving responses, code, or reasoning chains. A fundamentally more ambitious form of self-evolution targets the \emph{algorithms and architectures} themselves: rather than improving what the system produces, these methods improve \emph{how the system operates}. In this section, we analyze systems that use evolutionary computation, program search, and self-modifying code to discover novel algorithms, optimize agent architectures, and achieve open-ended improvement.

\subsection{LLMs as Evolutionary Operators}
\label{sec:llm-evolutionary}

The application of LLMs to evolutionary computation represents a paradigm shift. Traditional evolutionary algorithms rely on hand-designed mutation and crossover operators that make small, random changes to candidate solutions. LLMs, by contrast, can propose \emph{semantically meaningful} modifications informed by understanding of the problem domain.

\subsubsection{OPRO: LLMs as Optimizers}
\label{sec:opro}

OPRO (Optimization by PROmpting)~\citep{yang2023opro} demonstrates that LLMs can serve as general-purpose optimizers when the optimization problem is expressed in natural language. Given a history of (solution, score) pairs, the LLM proposes new candidate solutions expected to achieve higher scores.

\paragraph{Method.}
OPRO formulates optimization as a meta-prompting problem:
\begin{enumerate}
    \item Present the LLM with a prompt containing the objective function description and a history of previous solutions with their scores.
    \item The LLM generates new candidate solutions.
    \item Candidates are evaluated using the true objective function.
    \item The (solution, score) pairs are appended to the history.
    \item Repeat until convergence or budget exhaustion.
\end{enumerate}

\paragraph{Significance.}
While OPRO was initially applied to mathematical optimization (linear regression, traveling salesman), its framework generalizes naturally to \emph{any} domain where candidates can be expressed as text and evaluated by an automated scorer. This text-in, text-out optimization paradigm underlies virtually all LLM-driven evolutionary systems that follow.

\subsubsection{AlphaEvolve: Gemini-Powered Algorithm Discovery}
\label{sec:alphaevolve}

AlphaEvolve~\citep{novikov2025alphaevolve} represents the state of the art in LLM-driven evolutionary algorithm discovery. Developed at Google DeepMind, it combines dual LLM models with MAP-Elites quality-diversity search to discover and optimize algorithms across mathematics, engineering, and computer science.

\paragraph{Architecture.}
AlphaEvolve employs a multi-component architecture:
\begin{itemize}
    \item \textbf{LLM Ensemble}: Gemini Flash for fast, broad exploration and Gemini Pro for deep, targeted improvements. This dual-model strategy balances exploration (many diverse ideas) with exploitation (key refinements).
    \item \textbf{Evolutionary Database}: Maintains a growing population of candidate programs with associated fitness scores and behavioral descriptors.
    \item \textbf{Automated Evaluators}: Domain-specific objective functions that score candidate programs.
    \item \textbf{MAP-Elites}: A quality-diversity algorithm that maintains a map of high-performing solutions across behavioral dimensions.
\end{itemize}

\paragraph{Evolution loop.}
The AlphaEvolve loop operates as follows:
\begin{enumerate}
    \item \textbf{Selection}: Sample parent programs from the archive, weighted by fitness and novelty.
    \item \textbf{Variation}: The LLM proposes mutations or crossovers, producing child programs: $p_{\text{child}} = \text{LLM}(p_{\text{parent}}, \text{evolution\_context})$.
    \item \textbf{Evaluation}: Automated evaluators score the child programs on target objectives.
    \item \textbf{Archive Update}: Apply MAP-Elites update rule---the child replaces the current elite in its behavioral niche only if it achieves higher fitness: $\text{Archive}[b(p_{\text{child}})] = p_{\text{child}}$ if $f(p_{\text{child}}) > f(\text{Archive}[b(p_{\text{child}})])$.
\end{enumerate}

\paragraph{Results.}
AlphaEvolve's results are remarkable across multiple domains~\citep{novikov2025alphaevolve}:

\begin{itemize}
    \item \textbf{Mathematical Discovery}: Found an algorithm for $4\times4$ complex matrix multiplication using only \textbf{48 scalar multiplications}---the first improvement over Strassen's algorithm in 56 years. Across 50+ mathematical problems, AlphaEvolve re-discovers state-of-the-art solutions for 75\% and finds improvements for 20\%.
    \item \textbf{Infrastructure Optimization}: Recovered \textbf{0.7\% of worldwide Google compute} through improved data center scheduling (Borg system). Achieved a \textbf{23\% speedup} in FlashAttention kernel tiling and a \textbf{32\% speedup} in attention computation for LLM training.
\end{itemize}

\paragraph{Significance.}
AlphaEvolve demonstrates that LLM-driven evolution can produce results of genuine scientific and engineering significance, not merely incremental benchmark improvements. The combination of quality-diversity search with dual-model LLM variation is a template that future self-evolving systems are likely to follow.

\subsubsection{FunSearch: Mathematical Discovery through Program Evolution}
\label{sec:funsearch}

FunSearch~\citep{romeraparedes2023funsearch}, also from Google DeepMind, preceded AlphaEvolve and established the feasibility of LLM-driven mathematical discovery. FunSearch evolves entire programs (functions) rather than optimizing prompts or architectures.

\paragraph{Method.}
FunSearch operates similarly to AlphaEvolve but at the level of individual functions:
\begin{enumerate}
    \item Start with a seed function implementing a known algorithm.
    \item The LLM proposes modifications to the function body.
    \item An automated evaluator scores the modified function (correctness + performance).
    \item Best-performing functions are maintained in a population.
\end{enumerate}

\paragraph{Key result.}
FunSearch discovered new constructions for the cap set problem and the bin-packing problem, outperforming the best known human-designed algorithms~\citep{romeraparedes2023funsearch}. This demonstrated for the first time that LLM-driven program evolution could advance the frontier of mathematical knowledge.

\subsection{Agent Architecture Search}
\label{sec:agent-architecture-search}

A natural extension of evolutionary code discovery is the evolution of agent architectures themselves---automatically searching the space of possible agent designs to find configurations that outperform hand-crafted systems.

\subsubsection{ADAS: Automated Design of Agentic Systems}
\label{sec:adas-detail}

ADAS (Automated Design of Agentic Systems)~\citep{hu2024adas} defines the research area of automated agent design and introduces Meta Agent Search as its foundational method.

\paragraph{Core insight.}
The critical insight behind ADAS is that programming languages are Turing-complete, meaning the space of all valid Python programs that implement agents encompasses \emph{every possible agent architecture}. This includes novel prompting strategies, tool compositions, multi-step reasoning chains, recursive self-referential patterns, and any computable agent behavior.

\paragraph{Meta Agent Search.}
The Meta Agent Search algorithm operates as:
\begin{enumerate}
    \item A \textbf{Meta Agent} (an LLM) reviews previously discovered agent designs and their performance scores.
    \item The Meta Agent writes code for a new agent architecture: $a_{\text{new}} = \text{LLM}(D, \text{history}, \text{search\_strategy})$, where $D$ is the archive of discovered agents.
    \item The new agent is evaluated on benchmark tasks.
    \item The (code, score) pair is added to the archive $D$.
    \item The best agent is returned: $a^* = \arg\max_{a \in D} \text{score}(a)$.
\end{enumerate}

\paragraph{Results.}
Meta Agent Search discovered agents that outperform hand-designed state-of-the-art systems on ARC (Abstraction and Reasoning Corpus) and GFootball~\citep{hu2024adas}. Crucially, the discovered agents exhibit \textbf{cross-domain transfer}: agents designed for ARC transfer to GFootball, and agents discovered using GPT-3.5 perform well when executed with GPT-4. This transferability suggests that Meta Agent Search discovers general agent design principles rather than overfitting to specific tasks.

\paragraph{Recognition.}
ADAS received the \textbf{ICLR 2025 Outstanding Paper Award} and was recognized at the NeurIPS 2024 Open-World Agent Workshop, establishing agent architecture search as a major research direction.

\subsubsection{Darwin G\"odel Machine: Open-Ended Self-Improvement}
\label{sec:dgm-detail}

The Darwin G\"odel Machine (DGM)~\citep{zhang2025dgm} extends ADAS by combining it with open-ended evolution, creating a system that can improve itself without bound.

\paragraph{Architecture.}
DGM adds three key mechanisms on top of Meta Agent Search:
\begin{itemize}
    \item \textbf{Open-ended archive}: Unlike ADAS, which searches for a single best agent, DGM maintains a diverse archive $A = \{a_1, \ldots, a_n\}$ of agent implementations with fitness scores.
    \item \textbf{Diversity-biased sampling}: Parents are sampled from the archive with probability proportional to both fitness and diversity: $P(a_i) \propto f(a_i) \times \text{diversity\_bonus}(a_i, A)$.
    \item \textbf{Self-modifying code}: The LLM proposes modifications to the parent agent's source code: $a_{\text{child}} = \text{LLM}(a_{\text{parent}}.\text{code}, \text{edit\_instruction})$.
\end{itemize}

\paragraph{Results.}
DGM achieves dramatic improvements on software engineering benchmarks~\citep{zhang2025dgm}:
\begin{itemize}
    \item \textbf{SWE-bench}: 20.0\% $\rightarrow$ \textbf{50.0\%} Verified (a 2.5$\times$ improvement).
    \item \textbf{Polyglot} (multi-language coding): 14.2\% $\rightarrow$ \textbf{30.7\%} (a 2.2$\times$ improvement).
\end{itemize}
These results outperform all hand-designed agent architectures on the same benchmarks. Moreover, DGM discovered novel agent structures not imagined by human designers, demonstrating the creative potential of open-ended evolutionary search.

\paragraph{Philosophical connection.}
DGM is directly inspired by Schmidhuber's G\"odel Machine~\citep{schmidhuber2003godel}, a theoretical construct that can rewrite its own code in a provably optimal way. While DGM lacks the formal proof guarantees of the original G\"odel Machine, it operationalizes the same principle of self-referential code modification in a practical system.

\subsubsection{G\"odel Agent: Runtime Self-Modification via Monkey Patching}
\label{sec:godel-agent-detail}

G\"odel Agent~\citep{yin2025godelagent} takes a different approach to self-referential improvement: instead of maintaining an archive and evolving across generations, it modifies its own running code \emph{at runtime} through Python monkey patching.

\paragraph{Method.}
G\"odel Agent's self-modification loop:
\begin{enumerate}
    \item \textbf{Execution}: The agent executes on a batch of tasks, collecting results and feedback.
    \item \textbf{Self-analysis}: The agent reads its own source code and analyzes its performance: $A.\text{self\_analyze}(\text{code}, \text{results}, \text{feedback}) \rightarrow \text{improvement\_plan}$.
    \item \textbf{Self-modification}: The LLM proposes specific code changes (monkey patches) that modify the agent's methods at runtime: $A' = \text{LLM}(A.\text{code}, \text{feedback}, \text{objective})$.
    \item \textbf{Validation}: Modified agent runs on a validation set. Changes are accepted only if $\text{score}(A') > \text{score}(A)$; otherwise, a rollback occurs.
\end{enumerate}

\paragraph{Results.}
G\"odel Agent outperforms hand-designed agents across multiple benchmarks including HotPotQA (+8--15\%) and ALFWorld (significant improvement)~\citep{yin2025godelagent}. A key feature is continuous improvement: the agent's performance increases monotonically with task experience, and it discovers strategies not anticipated by human designers.

\paragraph{Comparison with DGM.}
G\"odel Agent and DGM represent two distinct approaches to the same goal:
\begin{itemize}
    \item \textbf{DGM} operates at the \emph{population level}---maintaining a diverse archive of agents and evolving across generations. It is more exploratory and better suited for discovering fundamentally novel architectures.
    \item \textbf{G\"odel Agent} operates at the \emph{individual level}---a single agent modifies itself at runtime. It is more focused and better suited for rapid, incremental improvement in deployment.
\end{itemize}
Both share the core insight that agent architectures should not be fixed at design time but should evolve through experience.

\subsection{Zero-Data Self-Play and Autonomous Curriculum Generation}
\label{sec:zero-data}

A critical bottleneck for self-evolving systems is the reliance on external data for training and evaluation. Several recent systems eliminate this dependency entirely by having agents generate their own training tasks through self-play.

\subsubsection{Absolute Zero: Self-Play Reasoning with Zero Data}
\label{sec:absolute-zero}

Absolute Zero~\citep{zhao2025absolutezero} introduces a paradigm where a single model simultaneously proposes tasks and solves them, using only self-generated data verified by a code executor.

\paragraph{Method.}
The Absolute Zero framework defines three roles, all played by the same model:
\begin{itemize}
    \item \textbf{Task Proposer}: Generates new reasoning tasks (code, math, logic problems).
    \item \textbf{Task Solver}: Attempts to solve the proposed tasks.
    \item \textbf{Code Executor}: Provides automated verification of solutions.
\end{itemize}

The RL training signal drives improvement in both roles:
\begin{align}
    t &= \text{LLM}_{\text{proposer}}(\text{curriculum\_state}) \\
    s &= \text{LLM}_{\text{solver}}(t) \\
    r &= \text{Execute}(s, t) \in \{0, 1\} \\
    R_{\text{proposer}} &\propto \text{difficulty}(t) \times \text{solvability}(t)
\end{align}

The proposer's reward is proportional to task difficulty (hard enough to be educational) and solvability (not trivially impossible), creating an automatic curriculum that maintains appropriate challenge.

\paragraph{Results.}
AZR-Coder-7B, trained with \textbf{zero external data}, achieves state-of-the-art results among 7B models on HumanEval+, MBPP+, and LiveCodeBench~\citep{zhao2025absolutezero}. Remarkably, the model acquires mathematical reasoning capability (competitive on GSM8K and MATH) despite being trained only on code self-play, demonstrating cross-domain transfer from code to mathematical reasoning.

\paragraph{Significance.}
The principle ``stronger models are better at making themselves stronger''~\citep{zhao2025absolutezero} has profound implications. It suggests a positive feedback loop: as models improve, their self-generated curricula become more effective, driving further improvement. This autocatalytic dynamic is a hallmark of true open-ended evolution.

\subsection{Open-Source Implementations}
\label{sec:opensource-evo}

The systems discussed above have spawned a growing ecosystem of open-source implementations that make evolutionary code discovery accessible to the broader research community.

\paragraph{OpenEvolve} is a community implementation of AlphaEvolve's concepts, providing an accessible framework for LLM-driven program evolution. It supports custom evaluation functions and has been applied to mathematical optimization, code generation, and algorithm discovery tasks.

\paragraph{SE-Agent}~\citep{seagent2025} applies trajectory-level evolution to software engineering, achieving 80\% on SWE-bench Verified (the top open-source result at time of publication). Its Revision-Recombination-Refinement mechanism operates across multiple agent trajectories, exchanging information between successful and failed attempts to drive cumulative improvement.

\paragraph{AgentEvolver}~\citep{agentevolver2025}, developed at Alibaba DAMO Academy, implements self-questioning, self-navigation, and self-attribution mechanisms that allow agents to autonomously identify improvement opportunities and modify their own behavior.

\subsection{Synthesis: From Self-Improvement to Self-Modification}
\label{sec:synthesis-evo}

The systems analyzed in this section reveal a clear progression in the depth of self-evolution:

\begin{enumerate}
    \item \textbf{Output improvement} (Section~\ref{sec:self-improvement}): The system improves what it produces---better responses, code, and reasoning---without changing its own structure.
    \item \textbf{Algorithm discovery} (Sections~\ref{sec:llm-evolutionary}--\ref{sec:opensource-evo}): The system discovers new algorithms and programs that outperform human-designed alternatives.
    \item \textbf{Architecture self-modification} (Sections~\ref{sec:agent-architecture-search}--\ref{sec:zero-data}): The system modifies its own architecture, discovering new ways to organize its components and processes.
\end{enumerate}

This progression mirrors biological evolution's trajectory from behavioral adaptation (learning within a lifetime) to genetic adaptation (evolution across generations). The key enabler at each level is the \textbf{evaluator}: a mechanism that can objectively score improvements. For output improvement, the evaluator is a reward model or test suite; for algorithm discovery, it is a mathematical verifier or performance benchmark; for architecture search, it is task-level performance on holdout sets.

The convergence of these three levels---output improvement, algorithm discovery, and architecture self-modification---represents the frontier of self-evolving AI. As we argue in Section~\ref{sec:open-problems}, the key challenge ahead is ensuring that these systems remain safe, controllable, and aligned with human values as they evolve beyond the scope of their original design.
